\section{The sixth problem among the twenty three}

The sixth problem is unusual in Hilbert's list because it asks for a program rather than a theorem, but it is connected with several of its neighbors. The second problem, the consistency of the axioms of arithmetic, is the logical counterpart of the sixth: Hilbert had reduced the consistency of the axioms of geometry to that of the axioms of arithmetic, and the sixth problem asks, among other things, that the axioms of a physical theory be shown to be compatible. G\"odel's incompleteness theorems of 1931~\cite{goedel1931} showed that a consistent formal system containing arithmetic cannot prove its own consistency, and the sixth problem inherited the limitation: an axiomatized physics can be no more certain than the mathematics in which it is written.

The nineteenth and twentieth problems concern the regularity of solutions of variational problems and the existence of solutions of boundary value problems. The nineteenth asked whether solutions of regular variational problems are always analytic. De Giorgi~\cite{degiorgi1957} in 1957 and Nash~\cite{nash1958} in 1958 answered it by proving H\"older continuity for solutions of elliptic equations with merely measurable coefficients, and Moser~\cite{moser1960} gave a new proof in 1960. De Giorgi's method was later used by Vasseur~\cite{vasseur2007} to give a new proof of the partial regularity theorem of Caffarelli, Kohn and Nirenberg~\cite{ckn1982} for the Navier Stokes equations, which Chapter~\ref{ch:boundary} discusses. The eighth problem, on the prime numbers, includes the Riemann hypothesis, which is the subject of the companion volume~\cite{polsonh8} and of Chapter~\ref{ch:open}. The thirteenth, on superpositions of functions, was answered for continuous functions by Kolmogorov and Arnold and appears in Chapter~\ref{ch:kolmogorov}.

\section{A worked example of the three readings}

The ideal gas shows the three readings in miniature.

\subsubsection*{The axiomatic reading} A monatomic ideal gas of $n$ molecules of mass $m_0$ in a box $\Omega$ has phase space $(\Omega\times\R^3)^n$ and Hamiltonian $H = \frac12m_0\sum_i|v_i|^2$, where $v_i$ is the velocity of molecule $i$. The microcanonical measure at energy $E$ is the normalized surface measure on $\{H = E\}$, and the canonical measure at inverse temperature $\beta$ is proportional to $e^{-\beta H}$. Both are probability measures in Kolmogorov's sense, and the equivalence of ensembles connects them: Theorem~\ref{thm:gibbs_conditioning} proves it for a finite alphabet, and Example~\ref{ex:maxwellian} carries it to the continuous gas.

\subsubsection*{The derivational reading} Under the canonical measure the velocities are independent Gaussian vectors with covariance $(k_BT/m_0)I_3$, the pressure $p$ on the walls is the mean momentum flux, and a short computation gives $p\,|\Omega| = nk_BT$. Under the microcanonical measure the same law holds in the limit $n\to\infty$ with $E/n$ fixed, by the Maxwell, Poincar\'e and Borel lemma (Theorem~\ref{thm:sphere_lemma} of Chapter~\ref{ch:poincare}). The macroscopic law is derived from the microscopic model by a law of large numbers.

\subsubsection*{The inferential reading} Suppose instead that we observe the velocities $v_1,\dots,v_k$ of a sample of $k$ molecules and wish to infer the temperature. With $\theta = k_BT/m_0$, the likelihood is $\prod_i(2\pi\theta)^{-3/2}e^{-|v_i|^2/2\theta}$. An inverse Gamma prior $\theta\sim\mathrm{IG}(s,b)$ is conjugate, and the posterior is
\[
\theta\mid v_{1:k}\sim\mathrm{IG}\Big(s+\frac{3k}{2},\;b+\frac12\sum_{i=1}^k|v_i|^2\Big),
\]
with posterior mean $(b+\frac12\sum|v_i|^2)/(s+\frac{3k}{2}-1)$, which tends to the equipartition estimate $\frac{1}{3k}\sum|v_i|^2$ as $k$ grows. Temperature, a macroscopic quantity, is a parameter of a de Finetti mixture over molecular velocities, and measuring it is Bayesian inference about that parameter.

Perrin's determination of Avogadro's number has the same structure one level up. Observing the displacements $\Delta_j$ of a colloidal particle over time intervals $\tau$, the model is $\Delta_j\sim N(0,2D\tau)$ independently in each coordinate, and Einstein's relation \eqref{eq:einstein_relation} gives $N_A = RT/(6\pi\eta aD)$. With a prior on $D$ the posterior for $N_A$ follows by change of variables. With $n$ displacements the relative posterior standard deviation of $D$ is of order $n^{-1/2}$, and it combines in quadrature with the relative uncertainty of the particle radius $a$, which Perrin measured separately, so which of the two dominates depends on $n$ and on the accuracy of the radius. Perrin reported about $6.4\times10^{23}$ per mole from the translational Brownian motion and about $6.8\times10^{23}$ from the height distribution of an emulsion in sedimentation equilibrium~\cite{perrin1926}, which lie within about six and thirteen percent of the modern value $6.02\times10^{23}$. The agreement of independent methods was Perrin's argument for the reality of molecules.

\section{Hilbert's later view}

In his address \emph{Naturerkennen und Logik}~\cite{hilbert1930}, given at K\"onigsberg on 8 September 1930, Hilbert returned to the relation between mathematics and physics. The address presents the axiomatic method as a general method for formulating scientific problems precisely, observes that mathematics often anticipates physics, as the theory of conic sections anticipated planetary orbits and Riemann's geometry anticipated general relativity, and insists that the laws of the world come from experience and are tested by experience alone. Geometry belongs to physics in this respect: whether rigid bodies exist, and how they can be placed relative to one another, is a matter of experience. The address ends with the sentence later engraved on Hilbert's tombstone, \emph{Wir m\"ussen wissen, wir werden wissen}: we must know, we shall know. In the days before, at a conference on the epistemology of the exact sciences held in the same city, G\"odel had announced his incompleteness theorem in a remark in the closing discussion~\cite{goedel1931b}. For the sixth problem, the relevant claim in the address is that the axioms of a physical theory come from experience and answer to it. Read inferentially, the choice of axioms becomes the choice of a prior and a likelihood, to be tested against data, and Hilbert's optimism becomes the Bayesian claim that learning from data is a mathematical process with guarantees.
\section*{Exercises}
\begin{enumerate}
\item In the ideal gas example, compute the posterior variance of $\theta$ and show that the posterior coefficient of variation is asymptotically $\sqrt{2/(3k)}$.
\item Show that for hard sphere collisions \eqref{eq:hs_collision} momentum and kinetic energy are conserved, and that for fixed $\omega$ the collision map $(v_i,v_j)\mapsto(v_i',v_j')$ is an involution with Jacobian of absolute value one.
\item Show that in the Boltzmann Grad scaling $N\varepsilon^{d-1} = \alpha$ the total volume fraction occupied by the spheres in a fixed box tends to zero, so the gas is dilute.
\item Let $X$ be a separable Hilbert space, let $(\Phi_t)$ be a strongly continuous semigroup of bounded linear operators on $X$, let the observation operators $\mathcal O_j$ in \eqref{eq:observation_model} be bounded and linear, let $\Pi = N(0,C)$ and let the noise variables be $N(0,\sigma^2I)$. Prove that the posterior is Gaussian and that its pushforward to trajectories is a Gaussian measure on $C([0,T];X)$. Find its mean and covariance when $X = \R^d$ and $\Phi_t = e^{-tA}$ for a symmetric positive definite matrix $A$, a discretization of the heat equation.
\item Consider a grain of radius $0.2$ micrometers in water at room temperature. Using \eqref{eq:einstein_relation}, estimate the diffusion coefficient and the root mean square displacement along one coordinate axis in thirty seconds.
\end{enumerate}
